\chapter{The Dalit Movement}

\section{Nationalism and the Backward/Dalit Question}

\begin{KeyConcept}
\begin{itemize}
\item \textbf{'India is a nation without nationalism'} --- nationalism is \textbf{psychological unity cutting across identities} (tribal, Dalit, woman, Brahmin), and that unity is still 'in the making'.
\item The Congress fought the British, but \textbf{Ambedkar, Phule and Periyar fought the Congress itself} --- their enemy was not the British but the Brahminical structure that monopolised resources (just as the communist's enemy is the landlord/industrialist, the women's enemy is patriarchy). So there is no single 'we'.
\item Nationalism is \textbf{dynamic, not static}: it appears in \textbf{disaster nationalism} (the Kerala floods --- everyone donates) and \textbf{sports nationalism} (India--Pakistan matches) --- \textbf{fission and fusion}: people fuse in some contexts and separate in others.
\end{itemize}
\end{KeyConcept}

\section{SNDP Movement (Shri Narayana Guru)}

\begin{KeyConcept}
\begin{itemize}
\item The \textbf{SNDP movement} (Shri Narayana Dharma Paripalana Yogam), led by \textbf{Shri Narayana Guru} in \textbf{Kerala}, is the last backward-class movement.
\item Narayana Guru belonged to the \textbf{Ezhava (Izhavas)} community --- traditionally \textbf{toddy tappers}, a low caste (alcoholism and non-vegetarianism being associated with the low castes). They were discriminated, outcasted and segregated --- denied temple entry, education and a common water source.
\item The Ezhavas were like the \textbf{Nadars of Tamil Nadu} and the \textbf{Idigas of Karnataka} (toddy-tapping communities).
\item Narayana Guru urged them to \textbf{abandon alcoholism, non-vegetarianism and obscene songs}, and to uplift themselves spiritually. The movement's objectives were \textbf{education for the Ezhavas, recruitment to government services, access to roads and temples, and spiritual uplift}.
\item It brought \textbf{transformative structural changes} --- upward social mobility and a shift in the traditional distribution of power.
\item \textbf{Critique}: instead of fighting the upper castes, Narayana Guru reformed the low castes themselves --- a \textbf{Sanskritizing} (reformative) strategy, not a real struggle.
\end{itemize}
\end{KeyConcept}

\section{Common Patterns and Critiques (Ghanshyam Shah and M. S. A. Rao)}

\begin{KeyConcept}
\begin{itemize}
\item \textbf{Ghanshyam Shah} --- the underlying pattern of all backward-class movements is \textbf{striving for social mobility *within* the system rather than fundamentally transforming it} (historically, \textbf{Sanskritization} was resorted to for mobility). They began as challenges to traditional Hindu values and Brahmanism, but \textbf{broke down when members developed political aspirations} (Shahu Maharaj's Kshatriya claim, the Justice Party's jobs, the Self-Respect League's Dravida Kazhagam).
\item \textbf{Phule's counter-narrative}: there is no 'Ram Rajya' (Gandhi) --- the only Rajya is \textbf{Bali Rajya} (Bali was the real king; he was killed by \textbf{Vaman, an avatar of Vishnu = a Brahmin}, so Brahmins have always been the outsiders who killed the original kings).
\item \textbf{Ambedkar's counter-narrative} (in '\textbf{Beef and Broken Men}'): the original Dalits were \textbf{Buddhists} who challenged the Brahmins and were outcasted --- the modern Dalits are those outcasted Buddhists.
\item These \textbf{genealogies/stories are constructed to challenge the established discourse and establish a new discourse} that can bring social mobility to the community (just as Yadavs trace descent to Yadu Maharaj).
\item \textbf{M. S. A. Rao} --- 'who is backward?': \textbf{backwardness is a social-political construct, and 'backward' is not homogeneous}. The Constitution speaks of the socially and educationally backward, but after OBC reservation \textbf{backward class got confused with backward caste} (each state has its own criteria). \textbf{Satish Deshpande}: within the backward there are the \textbf{'backward forward'} (socially forward, economically backward --- Maratha, Jat, Patidar) and the \textbf{'backward backward'}.
\item \textbf{M. S. A. Rao's four types of backward movement}: (1) led by \textbf{upper non-Brahmin castes} (Maratha, Jat --- reservation demand); (2) led by the \textbf{low castes} (Phule, Narayana Guru, the Self-Respect movement); (3) led by \textbf{untouchables against the upper and other backward castes} (the Ambedkarite movements); (4) \textbf{tribal movements} (separate).
\end{itemize}
\end{KeyConcept}

\begin{QuickRevision}
\begin{itemize}
\item SNDP (Shri Narayana Guru, Kerala): Ezhavas (toddy tappers, like Nadars/Idigas); stop drinking/non-veg; education + govt jobs + temple/road access + spiritual uplift; transformative change; Sanskritizing/reformative critique.
\item Ghanshyam Shah: mobility within (not transformation of) the system; movements broke down with political aspirations.
\item Counter-narratives: Phule (Bali Rajya; Vaman/Brahmin killed the king) + Ambedkar (Beef and Broken Men; Dalits = outcasted Buddhists) $\rightarrow$ new discourse for mobility.
\item MSA Rao: backwardness = social-political construct; backward-forward vs backward-backward (Deshpande); 4 types (upper non-Brahmin; low caste; untouchables; tribal).
\end{itemize}
\end{QuickRevision}

\section{Pre-Ambedkarite Dalit Movements}

\begin{KeyConcept}
\begin{itemize}
\item Before Ambedkar, \textbf{Dalits had no consciousness} --- they simply accepted their subordinate role; protest was only at the \textbf{local/regional level} (no all-India movement, no connectivity).
\item Important pre-Ambedkarite movements:
\item \textbf{1. The Satnami movement} of the Chamars of Chhattisgarh (1820), led by \textbf{Guru Ghasidas} (regional only).
\item \textbf{2. The Adi Dharma movement} in Punjab (1920s), led by \textbf{Babu Mangu Ram Mugowalia} --- 'to Punjab what Jyotiba Phule was to Maharashtra' (Punjab has the largest SC population).
\item \textbf{3. The socio-political mobilisation of the Jatavs of Agra} --- untouchable \textbf{shoemakers/leather-workers} who prospered under British industrialisation and urbanisation, used their wealth to \textbf{claim Kshatriya status} (Sanskritization) for political representation. \textbf{Owen Lynch} ('Politics of Untouchability'): the Jatavs have \textbf{two reference groups} --- they bank on their own community for \textbf{votes (numbers)}, but the candidates \textbf{emulate the elites} (elite emulation) for political seats.
\end{itemize}
\end{KeyConcept}

\begin{QuickRevision}
\begin{itemize}
\item Pre-Ambedkarite: no Dalit consciousness; regional protest only.
\item Satnami (Chamars, Chhattisgarh, 1820, Guru Ghasidas); Adi Dharma (Punjab, 1920s, Babu Mangu Ram Mugowalia); Jatavs of Agra (shoemakers $\rightarrow$ Kshatriya claim; Owen Lynch --- two reference groups, elite emulation).
\end{itemize}
\end{QuickRevision}

\section{The Ambedkarite Phase (Poona Pact, Yola Conference, Navayana)}

\begin{KeyConcept}
\begin{itemize}
\item \textbf{Ambedkar}: '\textbf{I was born Hindu, but I will never die Hindu}' --- having faced the consequences of untouchability; \textbf{Gandhi} replied that religion is not a house one changes at will.
\item The two landmark events of this phase: the \textbf{Poona Pact (1930)} and the \textbf{Yola Conference (1956)}.
\item \textbf{Poona Pact (1930)} --- Ambedkar demanded a \textbf{separate electorate} for Dalits; Gandhi opposed it. The division between Hindus and Dalits, until then only a sentiment, now became \textbf{official}. Gandhi believed the untouchables could be integrated through social reform; Ambedkar believed \textbf{Hinduism itself had to be reformed}.
\item \textbf{Yola Conference (1956)} --- Ambedkar declared that there was \textbf{no escape for the untouchables except conversion to Buddhism}, and converted that very day (a symbol to create confidence; 'you were originally Buddhists --- come back home').
\item He also \textbf{reformed Buddhism}: Buddhism had become \textbf{Brahminical} (Mahayana's deification of the Buddha vs the original Hinayana), so he added a \textbf{fourth wheel --- Navayana} (the new Buddhism).
\item Ambedkar's slogan: \textbf{'Educate, Agitate, Organize.'} Overall, he created \textbf{Dalit consciousness} by integrating Dalits into the new Buddhism.
\item \textbf{Why Buddhism (and not Islam/Sikhism/Jainism)?} Because Buddha was the \textbf{first to challenge Brahminism}. Ambedkar also \textbf{rejected the communists} --- they failed to differentiate the working class from the Dalits (in India the working class \emph{is} largely the Dalits, whose problem is \textbf{discrimination}, not just payment); he therefore formed the \textbf{Independent Labour Party} (for both workers and Dalits), not a communist party.
\end{itemize}
\end{KeyConcept}

\begin{QuickRevision}
\begin{itemize}
\item Ambedkar: 'born Hindu, will never die Hindu'; Gandhi: religion is not a house.
\item Poona Pact (1930): separate electorate (Ambedkar) vs integration (Gandhi); official Hindu/Dalit split.
\item Yola Conference (1956): conversion to Buddhism; Navayana (fourth wheel; against Brahminical/Mahayana deification).
\item Educate, Agitate, Organize; Ambedkarism (coined by Gail Omvedt = Phule + Marx + Ambedkar); rejected communists (discrimination vs payment); Independent Labour Party.
\end{itemize}
\end{QuickRevision}

\section{Post-Ambedkarite Movements (Dalit Panther, RPI, BSP)}

\begin{KeyConcept}
\begin{itemize}
\item The three post-Ambedkarite movements: the \textbf{Dalit Panther movement} (Maharashtra, college students, 1972 --- named after the Black Panther movement of the US), the \textbf{Republican Party of India} (1956), and the \textbf{Bahujan Samaj Party} (1984 --- rooted in \textbf{DS4/BAMCEF}).
\item All three aimed to \textbf{transform the system in favour of the Dalits}, running on Ambedkar's spirit of '\textbf{Educate, Agitate, Organize}'.
\item The \textbf{BSP} allied with the \textbf{Bahujans} (OBCs and tribals) to fight their common enemy --- the Brahmins.
\item The Dalit movement's four resources to claim: \textbf{education (reservation), temple entry (Mahad Satyagraha), political representation (RPI, BSP), and government jobs (the Independent Labour Party)}.
\end{itemize}
\end{KeyConcept}

\begin{QuickRevision}
\begin{itemize}
\item Post-Ambedkarite: Dalit Panther (1972, college students, from Black Panther), Republican Party of India (1956), BSP (1984, from DS4/BAMCEF).
\item All run on 'Educate, Agitate, Organize'; BSP allies Bahujans (OBC + tribal) vs Brahmins.
\item Four resources: education, temple entry, political representation, government jobs.
\end{itemize}
\end{QuickRevision}

\section{Dalit Assertion and the Contemporary Perspective}

\begin{KeyConcept}
\begin{itemize}
\item \textbf{Dalit assertion} = 'a claim of a community that is \textbf{dominant by its own right}, and does not depend on \textbf{Sanskritic rituals} for its upliftment.'
\item \textbf{Ghanshyam Shah's two types of Dalit movement}: (1) \textbf{Reformative} --- reforming the caste system (includes the \textbf{Bhakti movement, the Neo-Vedantic movement, and Sanskritization}); (2) \textbf{Alternative} --- creating an alternative socio-cultural structure (via \textbf{conversion} --- religious --- or via \textbf{education, wealth and power} --- secular). Both types use political means.
\item \textbf{Contemporary perspective}: in \textbf{Una (Gujarat)} a group of Dalits was assaulted by self-proclaimed \textbf{cow vigilants (gaurakshaks)} for skinning a dead cow --- recalling Ambedkar's story that the Brahmins were \textbf{originally beef-eaters} who turned vegetarian only when the Buddhists challenged them. Una became an \textbf{epicentre of anti-Brahminical assertion} (seen as a threat to the ruling party).
\item The \textbf{Rohith Vemula} incident and the \textbf{Hathras} case highlight that \textbf{untouchability has gone, but atrocity remains} --- and today atrocity against Dalits \textbf{provokes nationalist sentiments} (a \textbf{differential/nationalist mobilization}, unlike the pre-independence vertical/horizontal mobilizations).
\item Today the Dalit movement is \textbf{not limited to Dalit concerns} --- it speaks for \textbf{Muslims, farmers and women} (e.g. the Bhim Sena at the farmers' protest) --- reflecting the rise of \textbf{Bahujan nationalism}. The traditional \textbf{Brahmin--Dalit dichotomy has weakened with political democracy} (a Dalit President; a BSP with Brahmin members and funders).
\end{itemize}
\end{KeyConcept}

\begin{QuickRevision}
\begin{itemize}
\item Dalit assertion = dominant by its own right, no Sanskritic rituals.
\item Ghanshyam Shah: reformative (Bhakti, Neo-Vedantic, Sanskritization) vs alternative (conversion --- religious; education/wealth/power --- secular).
\item Contemporary: Una (gaurakshaks, skinning dead cow; Brahmins originally beef-eaters); Rohith Vemula, Hathras $\rightarrow$ 'untouchability gone, atrocity remains'; atrocity now provokes nationalist sentiments (differential/nationalist mobilization).
\item Dalit movement now speaks for Muslims/farmers/women (Bhim Sena) $\rightarrow$ Bahujan nationalism; Brahmin--Dalit dichotomy weakened (Dalit President; Brahmin members in BSP).
\end{itemize}
\end{QuickRevision}

